\section{Necessary conditions for the finite certificates}\label{sec:finite-criteria}
The universal dimension bound makes the classification a finite problem.
We record the necessary conditions used in its exact certificates. These
conditions concern arbitrary complex categorifications; positive unsquared
central dimensions are used only after an explicit pseudo-unitary reduction.

\subsection{Arithmetic and signed dimension characters}\label{finite:sec:finite-arithmetic}
\subsubsection{Conventions and standard obstructions}
A fusion ring has a distinguished basis $b_0=1,b_1,b_2,b_3$, nonnegative
integral structure constants, duality $i\mapsto i^*$, and Frobenius
reciprocity, with $N_{ij}^0=\delta_{i,j^*}$. We use the row convention
$(N_i)_{jk}=N_{ij}^k$. Commutativity is Lemma~\ref{core:lem:rank-four-algebra}.

Write $d_i=\FPdim(b_i)$ and $D=\sum_i d_i^2$. The four characters $\chi$
have formal codegrees
\[
f_\chi=\sum_i\chi(b_i)\chi(b_{i^*})
       =\sum_i|\chi(b_i)|^2.
\]
They are the eigenvalues, with multiplicity, of
\begin{equation}\label{finite:eq:casimir}
 B=\sum_i N_iN_{i^*}=\sum_iN_iN_i^{\mathsf T}.
\end{equation}
The following necessary conditions will be used in their stated scope.
The arithmetic inputs and permitted pseudo-unitary reductions are
Proposition~\ref{finite:prop:standard}.
The last assertion concerns a Galois conjugate; it does not claim
pseudo-unitarity of every categorification of the ring.

The Galois reduction used here is the one in
Section~\ref{sec:pivotal-bound}: an embedding of the dimension field
extends to an automorphism of $\CC$ and preserves all tensor and
spherical identities. For the two-dimension case we use precisely the
published statement in Proposition~\ref{finite:prop:standard}(5).

\subsubsection{Small arithmetic certificates}
For a nonzero algebraic integer $u$, put
$e_u=[\Q(u):\Q]$ and $n_u=|N_{\Q(u)/\Q}(u)|$.
\begin{lemma}[Normalized norm obstruction]\label{finite:lem:norm}
If $u/v$ is an algebraic integer, then
\begin{equation}\label{finite:eq:norm}
 \frac{v_p(n_u)}{e_u}\geq\frac{v_p(n_v)}{e_v}
 \qquad\text{for every rational prime }p.
\end{equation}
\end{lemma}
\begin{proof}
Taking the norm in $L=\Q(u,v)$ shows that
$n_u^{[L:\Q]/e_u}/n_v^{[L:\Q]/e_v}$ is an integer. Taking its
$p$-valuation gives the assertion.
\end{proof}
Thus a failed divisibility needs only two irreducible polynomials and one
prime. As a preliminary shortcut, a signed Casimir determinant not divisible
by a prime required by every possible $\Delta$ excludes the whole table:
if the minimal polynomial of $\Delta$ had norm divisible by that prime,
the characteristic polynomial of the Casimir would have constant term
divisible by it too.

\begin{lemma}[Prime certificates for fields]\label{finite:lem:field-prime}
Let $g\in\ZZ[t]$ be monic irreducible and
$K=\Q(\alpha)$, $g(\alpha)=0$.
\begin{enumerate}
\item If $K$ contains a root of a monic irreducible polynomial $f$, then
$\deg f$ divides $\deg g$.
\item Suppose $g$ splits into distinct linear factors modulo a prime $p$.
If $K$ contains all roots of a monic integral polynomial $f$, then $f$
splits into linear factors modulo $p$, with repetitions allowed.
\item If the squarefree reduction modulo $p$ of a monic irreducible
polynomial has irreducible factors of unequal degrees, its root field
cannot be cyclotomic.
\end{enumerate}
\end{lemma}
\begin{proof}
The first assertion is the tower formula. In the second, $p$ does not
divide the discriminant of $g$, hence not the index of
$\ZZ[\alpha]$ in the ring of integers of $K$. Every residue field
above $p$ is therefore $\mathbb F_p$. Reduce the linear factorization of
$f$ in the ring of integers of $K$ at any such prime.
For the third, Dedekind's factorization theorem identifies the factor
degrees with the cycle lengths of a Galois permutation of the roots.
A transitive abelian permutation group acts regularly after quotienting
its kernel, so all cycles of each element have the same length.
Every subfield of a cyclotomic extension is Galois and abelian.
\end{proof}
These are rejection certificates: failure to find a prime certificate
proves neither a field inclusion nor cyclotomicity. Positive root lists in
a number field are instead verified by multiplying their linear factors
and recovering the complete rational polynomial.

\subsubsection{Sphericalization and exhaustive signed tables}
Lemma~\ref{core:lem:family-induction} supplies the signed reduction.
For the finite search its coefficient conditions are
\begin{equation}\label{finite:eq:signed-bounds}
 |T_{ij}^{k}|\le N_{ij}^{k},\qquad
 T_{ij}^{k}\equiv N_{ij}^{k}\pmod2.
\end{equation}
All compatible signed tables, and all of their nonzero dimension
characters, are retained until a stated necessary condition excludes
them. A sign choice does not identify characters with different
ordinary or signed provenance.


The signed search is finite by \eqref{finite:eq:signed-bounds}. In the
self-dual case there are ten independent coefficients. The exhaustive
cover below retains all zero-denominator branches and all compatible
basis-sign transforms. A signed certificate rejects every remaining
table and every possible dimension-character factor.


\subsection{Central induction and two twist moments}\label{finite:sec:finite-induction}
Let $\C$ be pseudo-unitary, with its canonical spherical structure,
and let $F:\Z(\C)\to\C$ have right adjoint $I$.
For simples $X_i$ and $A$, write
$M_{iA}=[F(A):X_i]=[I(X_i):A]$. The four simple summands of $I(1)$
occur once, have twist one and dimensions $D/f_\chi$
\cite[Theorems 2.5 and 2.13]{Ostrik3}.
Let $L$ be their four restriction columns, including the central unit
column $(1,0,0,0)^{\mathsf T}$. Then
\begin{equation}\label{finite:eq:induction-gram}
 MM^{\mathsf T}=B,\qquad
 L\mathbf1=B_{\bullet0},\qquad
 B-LL^{\mathsf T}\text{ is a nonnegative integral Gram matrix}.
\end{equation}
Indeed $FI(X)=\bigoplus_j X_j\otimes X\otimes X_j^*$, and adjunction
gives the first identity. All remaining columns have first entry zero,
which gives the other assertions. In particular the remainder is positive
semidefinite and entrywise nonnegative.

\begin{lemma}[Gram saturation]\label{finite:lem:gram-saturation}
If $B=\sum_A m_Am_A^{\mathsf T}$ and some known columns already
satisfy $\sum_{A\text{ known}}(v^{\mathsf T}m_A)^2=v^{\mathsf T}Bv$,
then $v$ is orthogonal to every remaining column.
\end{lemma}
\begin{proof}
The sum of the remaining nonnegative squares is zero.
\end{proof}

\subsubsection{Higher vanishing indicators}
The following statement holds at arbitrary rank. For an object $V$, let
$I_{\rm low}(V)$ be the sum, with its multiplicities in $I(V)$, of those
simple constituents that occur in $I(1)$.

\begin{theorem}[Higher-indicator bound]\label{finite:thm:higher-indicator}
Let $\C$ be a pseudo-unitary fusion category over $\CC$, with its
canonical spherical structure, and let $m\geq1$. If
$\nu_1(V)=\cdots=\nu_m(V)=0$, then
\[
 (m+1)\FPdim(I_{\rm low}(V))\leq D\FPdim(V).
\]
More strongly, the same inequality holds with $I_{\rm low}(V)$ replaced
by the full twist-one part of $I(V)$. Equality forces every other
constituent to have a nontrivial $(m+1)$st root of unity as twist.
In particular, if every simple constituent of $V$ is non-self-dual, then
$3\FPdim(I_{\rm low}(V))\leq D\FPdim(V)$.
\end{theorem}
\begin{proof}
Write $w_A=[I(V):A]\FPdim(A)\geq0$ and $T=D\FPdim(V)$.
The indicator identity \cite[Theorem 4.1]{NS} is
$\sum_Aw_A\theta_A^k=D\nu_k(V)$. Negative moments vanish by complex
conjugation. Expanding the nonnegative polynomial
\[
 |1+z+\cdots+z^m|^2=(m+1)+
 \sum_{k=1}^m(m+1-k)(z^k+z^{-k})\qquad(|z|=1)
\]
gives $\sum_Aw_A|1+\theta_A+\cdots+\theta_A^m|^2=(m+1)T$.
Every twist-one constituent contributes $(m+1)^2w_A$; discarding the
remaining nonnegative terms proves the bound. Equality requires the
polynomial to vanish at every remaining twist, proving its stated support.
The last assertion follows from additivity and
$\nu_1(X)=\nu_2(X)=0$ for a non-self-dual simple $X$.
\end{proof}
For $m=1$ this is the low-part estimate explicit in the proof of
\cite[Theorem 2.21]{Ostrik3}. Its higher-indicator proof is included here. 

\subsubsection{The moment inequality}
For a nontrivial simple $X$, put $d=\FPdim(X)$ and $\epsilon=\nu_2(X)$.
The indicator formula \cite[Theorem 4.1]{NS} gives, with
$w_A=[I(X):A]\FPdim(A)\geq0$,
\begin{equation}\label{finite:eq:three-moments}
 \sum_A w_A=Dd,\qquad \sum_A w_A\theta_A=0,\qquad
 \sum_A w_A\theta_A^2=D\epsilon.
\end{equation}
Here the twists are roots of unity, $\epsilon\in\{-1,1\}$ for a self-dual
simple and $\epsilon=0$ otherwise.

\begin{proposition}[Two-moment bound]\label{finite:prop:optimized-moment}
If $L_X$ is any known contribution to the twist-one mass of $I(X)$, then
\[
 (3d+\epsilon)L_X\leq Dd(d+\epsilon).
\]
Consequently the following bounds require no choice of second indicator:
\begin{align}
 (3d+1)L_X&\leq Dd(d+1)&& (X\cong X^*),\label{finite:eq:optimized}\\
 3L_X&\leq Dd&& (X\not\cong X^*).\label{finite:eq:nonselfdual-third}
\end{align}
One may take $L_X=\FPdim(I_{\rm low}(X))=(LL^{\mathsf T}d)_X$.
\end{proposition}
\begin{proof}
For $P(z)=d+(d+\epsilon)z+dz^2$, expand $|P(z)|^2$ on the unit circle
and use \eqref{finite:eq:three-moments}. This gives
\[
 \sum_Aw_A|P(\theta_A)|^2=Dd(3d+\epsilon)(d+\epsilon).
\]
The twist-one terms contribute at least $(3d+\epsilon)^2L_X$.
Since $3d+\epsilon>0$, division gives the assertion.
The function $e\mapsto(d+e)/(3d+e)$ is increasing on $[-1,1]$;
replace $\epsilon$ by one in the self-dual case.
\end{proof}

The familiar first-moment bound $2L_X\leq Dd$ and the weaker Fej\'er
bound $9L_X\leq D(3d+2)$ follow as well. The checker uses
\eqref{finite:eq:optimized}--\eqref{finite:eq:nonselfdual-third}; these suffice for
all the low-induction inequalities needed here. Positivity is essential:
these bounds are not applied to arbitrary spherical categories with signed
dimensions.

\subsubsection{A finite restriction test}
\begin{proposition}[Exhaustive low-induction certificate]\label{finite:prop:low-certificate}
Fix the positive dimension character $d$. For each formal codegree $f$,
form the set of vectors $v\in\ZZ_{\geq0}^4$ satisfying
\begin{equation}\label{finite:eq:column-search}
 v_0=1,\qquad \sum_i v_id_i=D/f,\qquad
 v_i\leq\min\{B_{0i},\lfloor\sqrt{B_{ii}}\rfloor\}\quad(i>0).
\end{equation}
Reserve $(1,0,0,0)^{\mathsf T}$ for the central unit. Enumerate all choices
of the other three columns, with the codegree multiplicities retained,
whose sum together with the unit column is $B_{\bullet0}$.
If every such choice fails the entrywise or positive-semidefinite condition
in \eqref{finite:eq:induction-gram}, or one of
\eqref{finite:eq:optimized}--\eqref{finite:eq:nonselfdual-third}, then $R$ admits no
pseudo-unitary categorification.
\end{proposition}
\begin{proof}
A genuine restriction column satisfies its dimension equation, and
$M_{0A}=1$ on the four constituents of $I(1)$. The first row and the diagonal
of \eqref{finite:eq:induction-gram} give the two bounds in
\eqref{finite:eq:column-search}. Hence the actual low part occurs in the finite
list and satisfies every tested necessary condition.
\end{proof}

The computation is only linear algebra over a number field. Express the
four dimensions and $D/f$ in a rational power basis, row-reduce the dimension
equation, and enumerate its bounded free integer coordinates. When the
four dimensions are rationally independent, there is at most one possible
column for each $f$; a negative or nonintegral coefficient is already an
obstruction. When they are dependent, no uniqueness is presumed. Full
induction matrices, modular data and associators are never enumerated.




\subsection{Positive central moments and the sign object}
\label{finite:sec:frontier-signed-gauss}

We use the canonical spherical structure on the sphericalization
$\mathcal S$ of a hypothetical fusion category. Its realization as an
equivariantization for the canonical involutive pivotal functor gives
a central copy of $\operatorname{Rep}(C_2)$; see \cite[Section~5.3]{Bartlett}. Its central sign object
$\varepsilon$ has order two, dimension $-1$, and self-braiding $+1$.
Consequently its twist in $\Z(\mathcal S)$ is $-1$.
The two lifts of each original self-dual simple are self-dual; their
dimensions and second Frobenius--Schur indicators have opposite signs.
The latter assertion follows directly from the definition of the second
indicator: the two lifts have the same evaluation and coevaluation maps
after forgetting equivariance, while their pivotal maps differ by $-1$.

\begin{lemma}[A positive measure without pseudo-unitarity]
\label{finite:lem:positive-central-squares}
Let $\mathcal S$ be a spherical fusion category, and put
$d=\dim\mathcal S$. For $A\in\operatorname{Irr}\Z(\mathcal S)$,
let $a_A=\dim A$ and $\theta_A$ be its twist. Then $a_A$ is real and
\[
 \sum_A a_A^2=d^2,\qquad
 \sum_A a_A^2\theta_A=d,\qquad
 \sum_A a_A^2\theta_A^2=d\sum_V\dim(V)\nu_2(V).
\]
In particular the coefficients $a_A^2$ form a positive measure on roots
of unity, even when some dimensions of $\mathcal S$ are negative.
\end{lemma}
\begin{proof}
Simple dimensions in a spherical fusion category are totally real
\cite[Section~2.1]{Ostrik3}. The forgetful functor is spherical, so
$a_A$ is an integral linear combination of these real dimensions. Let $I$ be
central induction and $m_{V,A}=[I(V):A]$. Adjunction gives
$a_A=\sum_Vm_{V,A}\dim V$.
The indicator identity of \cite[Theorem~4.1]{NS}
\[
 \operatorname{Tr}(\theta^n\vert_{I(V)})=d\nu_n(V)
\]
therefore gives, after multiplication by $\dim V$ and summation over $V$,
\[\sum_A a_A^2\theta_A^n=d\sum_V\dim(V)\nu_n(V).\]
For $n=1$ only the tensor unit contributes, and for $n=2$ this is the
claimed formula. The zeroth moment is the standard identity
$\dim\Z(\mathcal S)=d^2$. The twist is a root of unity by the
Vafa theorem for the center of a spherical fusion category.
\end{proof}

\begin{proposition}[A global codegree moment bound]\label{finite:prop:spherical-codegree-moment}
Let a commutative fusion ring have a spherical categorification, with
simple dimensions $d_i$, global dimension $\Delta=\sum_i d_i^2$ and
formal codegrees $f_j$, counted with multiplicity. Put
\[
 w=\sum_j f_j^{-2},\qquad
 S=\sum_{i=i^*}|d_i|,\qquad h=\Delta^2w.
\]
At every real embedding of its dimension field one has
\begin{equation}\label{finite:eq:spherical-codegree-moment}
 2(\Delta-h)^2\leq
 (\Delta^2-h)(\Delta^2+\Delta S-2h).
\end{equation}
Equivalently,
$\Delta(1-3w)+(1-w)S+4w-2/\Delta\geq0$.
For a self-dual pseudo-unitary category,
$S=\sum_i\operatorname{FPdim}(X_i)$.
\end{proposition}
\begin{proof}
The constituents of $I(1)$ are distinct, have twist one and dimensions
$\Delta/f_j$. Remove their squared mass $h$ from the positive measure
of Lemma~\ref{finite:lem:positive-central-squares}. The remaining mass is
$\Delta^2-h$, its first moment is $\Delta-h$, and its second moment is
$\Delta s-h$, where $s=\sum_i d_i\nu_2(X_i)\leq S$.
The weighted sum of squared real parts of twists is
$(\Delta^2+\Delta s-2h)/2$. Cauchy--Schwarz gives the displayed
inequality with $s$ in place of $S$; increasing $s$ to $S$ preserves it
because the remaining mass is nonnegative. Expanding and dividing by
$\Delta^3>0$ gives the equivalent expression.
The argument applies to every Galois conjugate spherical structure.
The rational number $w$ can be computed without finding codegree roots:
$w=\operatorname{Tr}(B^{-2})$ for the Casimir matrix $B$.
\end{proof}

This form of the bound applies to every possible spherical dimension
character, including those outside the Frobenius--Perron orbit. A pseudo-unitary obstruction alone would not exclude such branches.

\begin{proposition}[Determined-sign moment obstruction]
\label{finite:prop:signed-translate-moment}
For a self-dual rank-four sphericalization use the notation
$E=\sum_i x_i^2$, $u=\sum f_i^{-2}$ and $v=\sum g_i^{-2}$ from
Section~\ref{sec:core}. Put $s=1+\sum_{i=1}^3\nu_i x_i$, where
$\nu_i=\nu_2(X_i^+)$. At every real embedding one has
\begin{gather}
 u+v\le2,\qquad E(u+v-1)\le s\le E,\label{finite:eq:sign-mass}\\
 (2-u-v)\left(1-u-v+\frac{s}{E}\right)
          \ge\left(u-\frac1E\right)^2.
 \label{finite:eq:determined-sign-moment}
\end{gather}
The same indicator assignment must work at every embedding.
\end{proposition}
\begin{proof}
The constituents of $I(1)$ have dimensions $E/f_i,E/g_i$ and twist
one. Their disjoint sign translates have opposite dimensions and
have twists $+1$ on the ordinary constituents and $-1$ on the signed
ones, by Lemma~\ref{core:lem:translate-twists}.
Thus their squared mass and signed first-moment contribution are
\[
 h=E^2(u+v),\qquad C=E^2(u-v).
\]
The first-moment identity for $I(\varepsilon)$ is automatic, since
both reciprocal-codegree sums are one. Removing these two families
from the positive squared-dimension measure leaves mass
$M=4E^2-2h$, first moment $2E-h-C$, and second moment $4Es-2h$.
Lemma~\ref{core:lem:measure-tools} gives the two mass bounds and
\[
 (2E-h-C)^2\le(4E^2-2h)(2Es+2E^2-2h).
\]
Substituting $h,C$ and dividing by $4E^4$ gives
\eqref{finite:eq:determined-sign-moment}. The indicators are rational
signs and remain fixed under Galois conjugation.
\end{proof}

\subsection{Joint arithmetic, parity and moment certificates}
\label{finite:sec:adjoint-parity}
\begin{theorem}[Joint signed certificate]\label{finite:cor:joint-signed-conditions}
A proposed signed quotient $T$ and dimension character $\chi$ must
satisfy all the following conditions. Its codegrees are $d$-numbers,
its character values are cyclotomic integers, and each ordinary or
signed codegree belongs to the character field and divides
$E=\sum_i\chi(b_i)^2$ integrally. Moreover there must be signs
$\nu_i\in\{1,-1\}$ such that
\begin{equation}\label{finite:eq:signed-parity}
 N_{j,j}^{i}-\nu_i T_{j,j}^{i}\equiv0\pmod4\qquad(1\le i\le3,\ 0\le j\le3),
\end{equation}
and \eqref{finite:eq:sign-mass}--\eqref{finite:eq:determined-sign-moment}
hold at every real embedding with $x_i=\chi(b_i)$ and that same
indicator assignment. Ordinary and signed occurrences of equal
codegrees retain their provenance and multiplicity.
\end{theorem}
\begin{proof}
The arithmetic statements are Lemma~\ref{core:lem:family-induction}
and the standard codegree and cyclotomicity inputs. The two lifts
of $b_i$ occur in an adjoint product with multiplicities
$(N_{j,j}^i+T_{j,j}^i)/2$ and $(N_{j,j}^i-T_{j,j}^i)/2$ and have
opposite indicators. By Lemma~\ref{core:lem:adjoint-parity}, the lift
with indicator $-1$ occurs evenly. This is exactly
\eqref{finite:eq:signed-parity}. These are the same indicators as
in Proposition~\ref{finite:prop:signed-translate-moment}; hence the
conditions must be imposed jointly, with one assignment at all
embeddings. There is no additional twist-sign search.
\end{proof}

For a spherical character the same parity observation improves
Proposition~\ref{finite:prop:spherical-codegree-moment}: replace its
upper bound $S$ by
\begin{equation}\label{finite:eq:parity-spherical-bound}
 S_{\rm par}=1+\sum_{i\in J}d_i+
       \sum_{\substack{i>0,\ i=i^*\\i\notin J}}|d_i|,
 \qquad J=\{i>0:i=i^*,\ N_{j,j^*}^{i}\text{ odd for some }j\}.
\end{equation}
Indeed parity forces $\nu_2(X_i)=1$ for $i\in J$, so
$s\le S_{\rm par}$. The remaining mass multiplying this upper bound
is nonnegative, preserving the inequality at each embedding.

\begin{lemma}[Parity-restricted signed enumeration]
\label{finite:cor:parity-restricted-enumeration}
Write the unsigned and signed tuples as $N=(N_0,\ldots,N_9)$ and
$p=(p_0,\ldots,p_9)$ in the convention \eqref{eq:tuple-convention}.
Let $J_1=\{0,3,5\}$, $J_2=\{1,6,8\}$ and $J_3=\{2,7,9\}$.
It suffices to search the union over $\nu\in\{1,-1\}^3$ of
\[
 |p_k|\le N_k,\qquad p_k\equiv\nu_iN_k\pmod4\ (k\in J_i),
 \qquad p_4\equiv N_4\pmod2.
\]
After identical and empty domains are removed, these are at most
eight pairwise disjoint boxes, also after the chosen sign normalization.
\end{lemma}
\begin{proof}
The congruences are \eqref{finite:eq:signed-parity}. For a fixed $i$,
the two residue choices agree when all $N_k$, $k\in J_i$, are even;
otherwise an odd coordinate makes them disjoint. Since the three
$J_i$ are disjoint, the product boxes have the same property.
A basis-sign change multiplies $T_{j,j}^i$ and $\nu_i$ by the same
sign, so their product is unchanged. The normalization therefore
preserves the union and loses no categorical possibility.
\end{proof}
Identical boxes do not identify their indicator assignments: every
compatible assignment remains available in the joint moment test.
The common bounded signed cover in
Section~\ref{finite:sec:shared-signed-cover} applies the universal
arithmetic conditions once, then tests all compatible transforms
against each ordinary ring.

\subsection{Uniform applications of the moment bounds}\label{finite:sec:applications}
The following consequences replace separate low-rank exclusion theorems.
Throughout, the two-dimension Galois reduction is used only in the form
stated in Proposition~\ref{finite:prop:standard}(5).

For the twist-sector observations in the next paragraph, work in a
pseudo-unitary category with its canonical spherical structure.
If $g$ is invertible of order two, its second indicator is
$\epsilon\in\{-1,1\}$. Since $\FPdim(g)=1$, equality in the triangle
inequality for the second indicator forces $\theta_A^2=\epsilon$ for
every constituent $A$ of $I(g)$. More generally, if $H_X$ is any known
mass in $I(X)$ supported on squared twist one, taking real parts of the
second moment gives
\begin{equation}\label{finite:eq:square-sector}
 2H_X\leq D(d_X+\nu_2(X))\leq D(d_X+1).
\end{equation}
Indeed the remaining real trace is at least the negative of its mass.
For a non-self-dual simple, equality in
$3W_X\leq Dd_X$ forces the complementary twists to be primitive cube
roots by Theorem~\ref{finite:thm:higher-indicator}.

\subsubsection{The rank-two subcategory obstruction}
\begin{proposition}\label{finite:prop:ranktwo}
If a fusion ring contains a based fusion subring with basis $1,x$ and $x^2=1+ax$, where $a\geq2$, it has no complex fusion categorification.
\end{proposition}
\begin{proof}
The indicated rank-two based subring would be the Grothendieck ring
of a fusion subcategory. Ostrik's rank-two classification
\cite[Theorem~1]{OstrikRankTwo} allows only $x^2=1$ and $x^2=1+x$.
\end{proof}
The criterion applies whenever such a rank-two based subring is present.

\subsubsection{A uniform diagonal two-invertible obstruction}
\begin{theorem}\label{finite:thm:diagonal}
Let $R_m$ have self-dual basis $1,g,X,Y$ and products
\[
g^2=1,\quad gX=Y,\quad gY=X,\qquad
X^2=Y^2=1+m(X+Y),\quad XY=g+m(X+Y).
\]
For every integer $m\geq3$, the ring $R_m$ has no fusion categorification over $\CC$.
\end{theorem}
\begin{proof}
Its positive dimensions are $(1,1,d,d)$, where $d=m+\sqrt{m^2+1}$ and $d^2=1+2md$. The two-dimension reduction allows a pseudo-unitary assumption. Put $D=4+4md$ and $V=X\oplus Y$.

The four dimensions in $I(1)$ are $1,d^2,1+md,1+md$. Moreover
\[
FI(1)=4\cdot1+2mX+2mY.
\]
In particular the low part of $I(g)$ is empty. Dimension comparison shows that the total multiplicities of $X,Y$ in the last three restrictions are $2m,m,m$. Therefore their known twist-one mass in $I(V)$ is, independently of its allocation between $X$ and $Y$,
\begin{equation}\label{finite:eq:diag-W}
W=2md^2+2m(1+md)=4m+6m^2d.
\end{equation}
The total mass in $I(V)$ is
\[
M=2Dd=8m+(16m^2+8)d.
\]

The Gram entry $B_{gg}=4$ implies that $I(g)$ is either twice one simple or a sum of four distinct simples. The first possibility contradicts its zero first twist trace. Write the four simple dimensions as $\beta_j=1+k_jd$: each restriction contains $g$ once, no unit, and $k_j$ copies of $X,Y$ in total. Thus $k_j\in\ZZ_{\geq0}$ and $\sum_jk_j=4m$.

Put $\epsilon=\nu_2(g)$. All four twists have square $\epsilon$. The zero first trace divides these four objects into the two possible twists, with mass $D/2$ at each. Since $d$ is irrational, each side has two objects and the two $k_j$ sum to $2m$. Each side therefore contributes at least
\[
2m+2m^2d
\]
to $I(V)$, by $k_1^2+k_2^2\geq(2m)^2/2$.

If $\epsilon=1$, the new twist-one mass is at least $2m+2m^2d$, disjoint from the low part. Adding \eqref{finite:eq:optimized} for $X,Y$ yields
\[
(3d+1)(6m+8m^2d)\leq 2Dd(d+1).
\]
After reduction by $d^2=1+2md$, the left side minus the right side is
\[
2(4m^2-m-4)+2(8m^3-4m^2-3m-4)d>0
\]
for $m\geq3$. For example, substituting $m=n+3$ gives coefficients
$8n^2+46n+58$ and $16n^3+136n^2+378n+334$, respectively.

Suppose instead $\epsilon=-1$. The four new twists are $\pm i$, and their total mass in $I(V)$ is at least
\[
H=4m+4m^2d.
\]
Remove this mass and the low mass $W$ from the positive measure of $I(V)$. For the remaining measure, the first cosine moment is $-W$, its total mass is at most
\[
M-W-H=(6m^2+8)d,
\]
and its second cosine moment is at most
\[
\frac{M+D(\nu_2(X)+\nu_2(Y))}{2}-W
\leq \frac M2+D-W=4+(2m^2+4m+4)d.
\]
Cauchy--Schwarz therefore requires
\[
W^2\leq (6m^2+8)d\bigl(4+(2m^2+4m+4)d\bigr).
\]
The difference between the two sides, reduced by the dimension equation, is
\begin{align*}
&8(m-2)(3m^3+3m^2+3m+2)\\
&\hspace{12mm}+8(m-2)(6m^4+6m^3+8m^2+5m+2)d>0.
\end{align*}
This excludes the second indicator choice as well.
\end{proof}
\subsubsection{Propagation from an invertible involution}
\begin{proposition}[Residual-involution criterion]\label{finite:prop:involution-propagation}
Let $\C$ be pseudo-unitary with commutative Grothendieck ring,
and let $g\ne1$ be an invertible
simple with $g^2=1$. For its actual low restriction matrix $L$, put
\[
 Q=B-LL^{\mathsf T},\qquad W_i=(LL^{\mathsf T}d)_i.
\]
Suppose $Q_{gg}=2$ and $0<W_g<D/2$. The two constituents of $I(g)$
outside $I(1)$ then have dimensions and twists
\[
 (D/2-W_g,1),\qquad(D/2,-1).
\]
If moreover $D/2-W_g=1$, their restriction columns are respectively
$e_g$ and $Q_{g\bullet}^{\mathsf T}-e_g$. Consequently, for every
nonunit simple $X_i\ne g$,
\begin{equation}\label{finite:eq:involution-propagation}
 2W_i+DQ_{gi}\leq D(d_i+\nu_2(X_i))\leq D(d_i+1).
\end{equation}
\end{proposition}
\begin{proof}
The residual Gram entry $Q_{gg}=2$ is a sum of squares of nonnegative
integer multiplicities. Thus exactly two additional simples occur,
each once. Equality in the triangle inequality for the second indicator
of $g$ forces all twists in $I(g)$ to have square $\nu_2(g)$.
The positive low mass has twist one, so $\nu_2(g)=1$.
The vanishing first trace divides the total dimension $D$ equally
between twists $1$ and $-1$. Since $0<W_g<D/2$, both additional
sectors have positive mass and each contains one of the two remaining
simples. Their dimensions follow.

When the first dimension is one, its restriction contains $g$ and has
Frobenius--Perron dimension one; hence its restriction is exactly $g$.
Write $v$ for the other restriction column. In the residual Gram product,
only these two columns contribute to row $g$, so
$Q_{gi}=\delta_{gi}+v_i$ for every $i$.
For $i\ne g$, the low constituents and the twist-minus-one constituent
therefore give disjoint squared-twist-one mass
$W_i+(D/2)Q_{gi}$ in $I(X_i)$. Equation~\eqref{finite:eq:square-sector}
proves the required bound.
\end{proof}
The criterion is finite and uniform: enumerate the admissible low parts
by Proposition~\ref{finite:prop:low-certificate} and check
\eqref{finite:eq:involution-propagation} whenever its hypotheses hold.
Uniqueness of an admissible low part is sufficient but is not assumed.
The automated dispatcher recognizes this obstruction from the complete
low-induction proof of any weakly integral input ring; it does not match
the tensor against a preselected exceptional ring.



\subsubsection{The non-self-dual quadratic family}

For an integer $m\geq0$, let $S(m)$ be the commutative based ring with basis
$\{1,g,x,y\}$, involution $g^*=g$, $x^*=y$, and multiplication
\[
 g^2=1,\qquad gx=y,\qquad gy=x,\qquad
 x^2=y^2=g+m(x+y),\qquad xy=1+m(x+y).
\]

\begin{theorem}\label{finite:thm:nonselfdual-quadratic}
For every $m\geq2$, the fusion ring $S(m)$ admits no fusion categorification
over $\CC$.
\end{theorem}

\begin{proof}
Put $d=m+\sqrt{m^2+1}$. Thus $d^2=1+2md$, and the Frobenius--Perron
dimensions are $(1,1,d,d)$. By Schopieray's Galois-conjugacy theorem for
two-dimension fusion rings \cite[Lemma~3.1]{Schopieray}, any fusion
categorification would have a pseudo-unitary Galois conjugate with the same
fusion rules. It is consequently sufficient to exclude pseudo-unitary
categorifications.

Suppose one exists, and use its canonical spherical structure. Its dimension is
$D=2+2d^2=4+4md$. The four characters of $S(m)$, displayed by columns, are
\[
 \begin{pmatrix}
 1&1&1&1\\
 1&1&-1&-1\\
 d&-d^{-1}&i&-i\\
 d&-d^{-1}&-i&i
 \end{pmatrix}.
\]
Their formal codegrees are $D$, $2+2d^{-2}$, $4$, $4$, respectively.
The formal-codegree description of induction therefore gives
\[
 I(\mathbf1)\cong\mathbf1\oplus A\oplus B\oplus C,
 \qquad
 \operatorname{FPdim}(A)=d^2=1+2md,
 \quad
 \operatorname{FPdim}(B)=\operatorname{FPdim}(C)=1+md,
\]
where the four central simples are pairwise nonisomorphic and occur with
multiplicity one. Indeed, every irreducible representation of this commutative
rank-four ring is one-dimensional. Furthermore,
\[
 F I(\mathbf1)\cong
 \bigoplus_{T\in\operatorname{Irr}\C}T\otimes T^*
 \cong 4\mathbf1\oplus2mX\oplus2mY,
\]
where $X^*\cong Y$ correspond to $x,y$. By adjunction each of $F(A)$,
$F(B)$, $F(C)$ contains $\mathbf1$ once, and the displayed decomposition
shows that none contains $g$. Since $d$ is irrational, comparison of dimensions
now forces
\[
 [F(A):X]+[F(A):Y]=2m,\qquad
 [F(B):X]+[F(B):Y]=[F(C):X]+[F(C):Y]=m.
\]
These statements determine the sum of the two low-part dimensions without
requiring the individual restriction multiplicities:
\begin{align*}
 L&=\operatorname{FPdim}(I_{\mathrm{low}}(X))+
    \operatorname{FPdim}(I_{\mathrm{low}}(Y))\\
  &=2m(1+2md)+2m(1+md)=4m+6m^2d.
\end{align*}
Both $X$ and $Y$ are non-self-dual. Applying
\eqref{finite:eq:nonselfdual-third} to each of them and adding gives $3L\leq2dD$.
But
\begin{align*}
 3L-2dD
 &=12m+18m^2d-2d(4+4md)\\
 &=12m+18m^2d-8d-8m(1+2md)\\
 &=4m+(2m^2-8)d>0
\end{align*}
for $m\geq2$, a contradiction.
\end{proof}



\subsection{A quadratic root-of-unity obstruction}\label{finite:sec:quadratic-family}

Let $R_n$ be the ordinary table in
\eqref{quad:eq:rational-three-normal-form}, with $m=n\ge1$ and
its tuple specified in the same calculation. Its positive dimension
row is $(1,\delta-1,\delta,\delta+1)$, where
\[
 \delta=\frac{3n+\sqrt{9n^2+4}}2,
 \qquad \delta^2=3n\delta+1,\qquad D=3+3\delta^2=6+9n\delta.
\]

\begin{lemma}\label{finite:lem:odd-quadratic-length}
Let $d>1$ be squarefree and let $f$ be the conductor of $\Q(\sqrt d)$:
$f=d$ if $d\equiv1\pmod4$ and $f=4d$ otherwise. Put
$\kappa=\sqrt{f/d}$ and $\lambda=(\varphi(f)-|\mu(f)|)/\kappa$.
Here $\varphi$ and $\mu$ denote Euler's totient and the M\"obius function.
If a sum of $k$ roots of unity equals $a+b\sqrt d$, with
$a,b\in\Q$, then $k\geq-a-\lambda b$.
\end{lemma}
\begin{proof}
Let $K=\Q(\sqrt d)$, let $L$ be a common cyclotomic overfield, and
write $T=\operatorname{Tr}_{L/K}/[L:K]$. Define
$\ell(a+b\sqrt d)=-a-\lambda b$. We prove $\ell(T(\xi))\leq1$
for every root of unity $\xi$.

If $m$ is the exact order of $\xi$ and $f\nmid m$, then
$K\cap\Q(\xi)=\Q$, because the conductor of $K$ is $f$.
Consequently $T(\xi)=\mu(m)/\varphi(m)$, whose absolute value is at most one.
If $m=fh$, let $\chi$ be the primitive quadratic character modulo $f$ and
write $\xi=\zeta_m^r$, with $(r,m)=1$. The induced quadratic Gauss sum gives
\[
 T(\xi)=
 \begin{cases}
 \displaystyle\frac{\mu(f)\mu(h)+\chi(r)\mu(h)\chi(h)\kappa\sqrt d}
                      {\varphi(f)\varphi(h)},&(f,h)=1,\\[6pt]
 0,&(f,h)>1.
 \end{cases}
\]
For clarity, this follows by splitting the units modulo $m$ into the two
values of $\chi$: their sum over $\chi=1$ is half the sum of the Ramanujan
sum $\mu(m)$ and the induced Gauss sum. The latter equals
$\chi(r)\mu(h)\chi(h)\sqrt f$ when $(f,h)=1$ and vanishes otherwise;
an additional power of a prime dividing $f$ cancels upon summing over its
last prime-power layer. The same formula includes nonsquarefree $h$ through
$\mu(h)=0$.
In the nonzero case, the definition of $\lambda$ gives
\[
 |\ell(T(\xi))|
 \leq\frac{|\mu(f)|+\lambda\kappa}{\varphi(f)\varphi(h)}\leq1.
\]
Applying $\ell\circ T$ to the sum proves the lemma.
\end{proof}

\subsubsection{The conductor-forty equality case}
\begin{lemma}\label{finite:lem:rootlength40}
If $k$ roots of unity have sum $a+b\sqrt{10}$, then $k\ge-a-8b$.
Equality forces each root to be either $-1$ or a primitive fortieth
root $\zeta_{40}^r$ with $\chi_{40}(r)=-1$. In particular, if eighteen
roots have sum $-2-2\sqrt{10}$, exactly two equal $-1$ and the
normalized absolute trace of the sum of their squares is $2$.
\end{lemma}
\begin{proof}
In Lemma~\ref{finite:lem:odd-quadratic-length}, $f=40$, $\kappa=2$ and
$\lambda=8$. For a root of order $m$ not divisible by $40$, equality
in the normalized trace bound requires
$-\mu(m)/\varphi(m)=1$, hence $m=2$.
For $m=40h$, the nonzero trace formula requires $(h,40)=1$, and
equality requires $\varphi(h)=1$. Thus $h=1$, and equality with the
positive sign means $\chi_{40}(r)=-1$. Equality in a sum requires
these conditions for every summand.
Primitive fortieth roots have normalized absolute trace zero, whereas
$-1$ has trace $-1$. The trace of $-2-2\sqrt{10}$ is $-2$, so precisely
two roots equal $-1$. Squaring a primitive fortieth root gives a
primitive twentieth root, also of trace zero since $\mu(20)=0$.
The other two squares are $1$, as asserted.
\end{proof}


\begin{proposition}\label{finite:prop:quadratic-family-n3-n4}
Let $n\geq2$, write $9n^2+4=t^2d$ with $d$ squarefree, and define
$\lambda$ as in Lemma~\ref{finite:lem:odd-quadratic-length}.
If $n=2$, or if
\begin{equation}\label{finite:eq:rootlength-parameter}
 n^2+nt\lambda>8n^2+4,
\end{equation}
then $R_n$ admits no spherical categorification over $\CC$.
The condition is an exact scalar arithmetic test for each parameter.
\end{proposition}
\begin{proof}
We first derive a necessary root-of-unity identity for a pseudo-unitary
categorification of $R_n$. Give such a category its canonical spherical
structure, and use $F,I$ for forgetful functor and central induction.
The four characters, displayed by rows, are
\[
 \begin{pmatrix}
 1&\delta-1&\delta&\delta+1\\
 1&\delta'-1&\delta'&\delta'+1\\
 1&0&1&-1\\
 1&1&-1&0
 \end{pmatrix},\qquad \delta'=-\delta^{-1}.
\]
The formal codegrees are $D,D'=3+3(\delta')^2,3,3$. Accordingly,
\[
 I(1)=1\oplus A\oplus B\oplus C,
 \qquad d_A=D/D'=1+3n\delta,\quad d_B=d_C=D/3=2+3n\delta.
\]
Every summand has multiplicity one and twist one. Comparing the coefficients
of $1$ and $\delta$ in their restriction dimensions shows that their columns are
\[
 (1,a,3n-2a,a),\quad
 (1,b,3n-1-2b,b+1),\quad
 (1,c,3n-1-2c,c+1).
\]
The first row of $MM^{\mathsf T}=Q:=\sum_U N_U^2$, for the full induction
matrix $M$, gives $a+b+c=3n-1$.

Here the low part is forced without searching over $a,b,c$. Put
$v=(0,1,-2,1)^{\mathsf T}$. Multiplication of the displayed matrices gives
$v^{\mathsf T}Qv=18$. On the three low columns, the values of $v$ are
\[
 6(a-n),\quad 6(b-n)+3,\quad6(c-n)+3.
\]
Their squares sum to at most $18$. The last two are odd multiples of three,
so their squares already sum to at least $18$. It follows that
$a=n$ and $\{b,c\}=\{n-1,n\}$. Thus
\[
 F(A)=(1,n,n,n),\quad F(B)=(1,n-1,n+1,n),\quad
 F(C)=(1,n,n-1,n+1).
\]
Lemma~\ref{finite:lem:gram-saturation} now makes $v$ vanish on every
remaining column. Such a column is therefore
\[
 (0,b-t,b,b+t),\qquad b\geq|t|,\quad b>0,
\]
and its central dimension is $3b\delta+2t$. Subtracting the four low columns
from $Q$ and comparing its entries gives
\begin{equation}\label{finite:eq:quadratic-family-gram}
 \sum b^2=6n^2+2,\qquad \sum bt=2n,\qquad\sum t^2=2.
\end{equation}
Thus precisely two high columns have nonzero $t$, which equals $1$ or $-1$.
Set $T=3n\delta+2$. Subtracting the first twist-trace equation for $X$ from
that for $Z$, then dividing by two, gives
\begin{equation}\label{finite:eq:quadratic-family-ttrace}
 \sum t(3b\delta+2t)\theta=-T.
\end{equation}
If the two signs of $t$ differ, their $b$ values differ by $2n$, by
\eqref{finite:eq:quadratic-family-gram}. Their positive dimensions then differ by
$6n\delta+4=2T$, contradicting the reverse triangle inequality in
\eqref{finite:eq:quadratic-family-ttrace}. Both signs cannot be negative. Hence both
are positive; write their $b$ values as $p,q$, with $p+q=2n$.
Their dimensions $\alpha=3p\delta+2$, $\beta=3q\delta+2$ satisfy
\[
 \alpha+\beta=2T,\qquad
 \alpha\beta=(9pq+4)\delta^2,\qquad T^2=(9n^2+4)\delta^2.
\]
If their twists are $\zeta,\eta$, equation
\eqref{finite:eq:quadratic-family-ttrace} implies
\[
 \zeta\eta^{-1}+\eta\zeta^{-1}
 =2-\frac{3(9n^2+4)}{9pq+4}.
\]
This rational number is an algebraic integer and lies in $[-2,2]$.
Since $pq\leq n^2$, it is at most $-1$. The value $-2$ would require
$36pq=27n^2-4$, which is impossible modulo nine. The value $-1$ forces
$pq=n^2$, hence $p=q=n$. The two twists then sum to $-1$ and are the
primitive cube roots of unity.

All other columns have $t=0$ and are of the form $(0,b,b,b)$. Their squared
multiplicities sum to $4n^2+2$. Write $S=\sum b^2\theta$ over these uniform
columns; it is a sum of $4n^2+2$ roots of unity, with repetitions.
The first twist-trace equation for $X$ gives
\[
 S=n^2-n\delta.
\]
For completeness, the low contribution is
$n(1+3n\delta)+(2n-1)T$; the two exceptional high columns contribute
$-(n-1)T$, so the equation is
$0=n(1+3n\delta)+nT+3\delta S$, which is equivalent to the claimed identity
using $\delta^{-1}=\delta-3n$.

Writing $9n^2+4=t^2d$ with $d$ squarefree, we have
\[
 S=-\frac{n^2}{2}-\frac{nt}{2}\sqrt d.
\]
Lemma~\ref{finite:lem:odd-quadratic-length} therefore requires
\[
 4n^2+2\geq\frac{n^2+nt\lambda}{2}.
\]
Condition~\eqref{finite:eq:rootlength-parameter} contradicts this inequality.
For $n=2$, the same reduction gives eighteen roots whose sum is
$S_1=-2-2\sqrt{10}$. The second indicator of $X$ is
$\epsilon\in\{1,-1\}$. The same twist-trace equation with twists squared,
using that the two exceptional twists still sum to $-1$, gives
\[
 S_2=S_1+\epsilon\frac{D}{3\delta}
     =-2+(2\epsilon-2)\sqrt{10}.
\]
For $\epsilon=-1$, Lemma~\ref{finite:lem:rootlength40} applied to the eighteen
squared roots requires at least $34$ roots. For $\epsilon=1$, the same
lemma applied to $S_1$ shows that the normalized absolute trace of $S_2$
is $2$, whereas $S_2=-2$. Both cases are impossible.

Finally, a spherical categorification has a dimension character with no zero entries,
so the displayed character table restricts it to the two quadratic characters.
These are Galois conjugate; conjugating its structure constants if necessary
would give a pseudo-unitary categorification. This completes the proof.
\end{proof}


